\chapter{Methodology}
\label{ch:methodology}

% The scientific-correctness backbone. Everything in Sec. "a-priori locks" must
% be fixed BEFORE results are read (no HARKing, no p-hacking; \cite{kerr1998}).

\section{Pipeline Under Evaluation}
\label{sec:method-architecture}

% TODO: brief -- the design is in Ch.~\ref{ch:system-design}. State precisely what
% is measured: the reading stage, the structuring stage, and the deterministic
% post-processing between them; plus the single-shot variant used as the
% architectural comparison arm.

\section{Corpus}
\label{sec:method-data}

% TODO: the German e-invoicing reference corpus \cite{zugferdcorpus} plus
% programmatically generated invoices; page rendering settings; how document variety
% (single/multi rate, credit notes, reverse charge, discounts, multi-page) is covered.
% Document the corpus per the datasheet convention \cite{gebru2018datasheets}.
% Be explicit that these are born-digital renders, not scans -- the limitation is
% discussed in Ch.~\ref{ch:limitations-future-work}.

\section{Ground Truth from Embedded Structured Data}
\label{sec:method-gt}

% TODO: the central methodological device. Hybrid e-invoices carry both a rendered
% page and an embedded machine-readable representation of the same content, so the
% reference data is obtained without manual labelling: render the page, hide the
% embedded data from the model, extract it independently as the reference.
% Cover: the extraction path, cross-validation of the reference with a second
% independent tool, and the \ac{EN}~16931 business-term mapping.

\section{Target Schema and Field Registry}
\label{sec:method-schema}

% TODO: the evaluated field set and how it is declared once, in one registry, that
% simultaneously drives generation prompts, scoring, and reporting -- so the schema
% cannot drift between the three. Cover flat fields and the repeating groups
% (line items, tax breakdown, payment discounts).

\section{Compliance-Aware Evaluation Suite}
\label{sec:method-metrics}

% TODO (fixed before results were read): per-field precision / recall / \ac{F1},
% micro and macro; field weighting by legal importance under \S14~\ac{UStG} and
% \S33~\ac{UStDV}; document-level compliance pass rate; input-tax-deductibility
% decision accuracy; per-field error analysis.
% State plainly which metric is the headline and why.

\section{Measurement Validity: Validating the Ruler}
\label{sec:method-measurement-validity}

% TODO: a distinctive part of this thesis's method, and worth several pages.
% The problem: the model sees the page while the reference stores canonical values,
% so a naive string comparison punishes correct answers. Examples of the classes
% found and corrected:
%   - controlled-vocabulary codes that never appear in print (the page says
%     "Umsatzsteuer", the reference stores a category letter);
%   - values printed together with their label ("Invoice no. 471102" vs "471102");
%   - optional totals that the reference omits but a model may emit as zero;
%   - document-level scalars that are ill-defined for some documents (a single
%     "tax rate" on an invoice carrying two different rates);
%   - content the page renders as free text while the reference stores it
%     structurally, or the reverse.
% Method: each correction is a separate, reviewed change with a recorded rationale
% and a regression test; corrections may only change REPRESENTATION handling and may
% never let a genuinely wrong value score as correct. State that discipline
% explicitly -- it is what separates fixing the ruler from inflating the score.
% Reference the decision records for the full chain; summarise the taxonomy here.

\section{Sealed Train and Validation Splits}
\label{sec:method-splits}

% TODO (fixed before results were read): a stratified split sealed before any
% adaptation work, with fingerprints committed so the partition is auditable after
% the fact; what stratification variables were used; the guards that keep validation
% documents out of prompt design and adaptation.

\section{Reading-Quality Proxy}
\label{sec:method-answerability}

% TODO: to compare candidate readers without running the whole pipeline for each,
% reading quality is measured directly: for each document, what share of the
% reference values is recoverable from the transcript (allowing for the usual
% German formatting variants of dates, decimal separators and account numbers).
% Report the measured correlation with the end-to-end metric, and be explicit that
% this is a proxy with a known blind spot -- a value can be present as a string yet
% unusable if the surrounding structure is destroyed.

\section{Error Attribution Design}
\label{sec:method-attribution}

% TODO: the three instruments, and why three are needed:
%   1. re-scoring the SAME stored predictions per field group, to see where the loss
%      is concentrated without re-running any model;
%   2. classifying every miss by whether the reference value was findable in the
%      transcript at all, which separates a reading failure from a structuring one;
%   3. re-running the identical structuring pass on reference-quality transcripts,
%      which yields an upper bound on the structuring stage independent of reading.
% State that all three arms use matched decoding settings, and that instrument 3 is
% the decisive one because it removes reading from the equation entirely.

\section{Pre-Registration and Honest Reporting}
\label{sec:method-hypotheses}

% TODO: the hypotheses, their registration dates, and the discipline: predictions
% recorded before the corresponding data was collected \cite{kerr1998}; each
% experiment cites the hypothesis it tests or is labelled exploratory; a refuted
% hypothesis is reported as such. Also state, here rather than late in the thesis,
% that some registered hypotheses were not evaluated after the scope was narrowed,
% and point to Ch.~\ref{ch:limitations-future-work}. The full register is in the
% appendix.

\section{Reproducibility and Hardware}
\label{sec:method-repro}

% TODO: configuration-as-data (every experiment knob in a versioned configuration
% file, validated before any model loads), deterministic decoding, seeding,
% experiment tracking, and the exact local hardware envelope; note where a rented
% GPU was used and why (one candidate exceeded the local memory envelope).
